\section{Extensions}
\label{sec-extensions}
\subsection{Unbalanced OT}
\label{sec-unbalanced}
\paragraph{Relaxed formulation.}
\eql{\label{eq-unbalanced-primal}
\UW_c(\al,\be)= \uinf{\pi \in \Mm_+(\Xx \times \Yy)} \int_{\Xx \times \Yy} c(x,y) \d\pi(x,y)
+ \Dd_{\psi_1}(\pi_1|\al) + \Dd_{\psi_2}(\pi_2|\be)
}
\(\UW_{c,\tau}(\alpha,\beta) = \inf_{\pi\geq0} \int c\d\pi + \tau\Dd_{\bar\psi_1}(\pi_1|\alpha) + \tau\Dd_{\bar\psi_2}(\pi_2|\beta).\)
\begin{prop}[Small-transport-scale limit for KL penalties]\label{prop-unbalanced-kl-hellinger}
Assume $c\geq0$ and $c(x,x)=0$ on a common space, and take $\bar\psi_1=\bar\psi_2=\KL$. Then
\[
\lim_{\tau\downarrow0}\frac{1}{\tau}\UW_{c,\tau}(\alpha,\beta)
=
\inf_{\rho\in\Mm_+(\Xx)} \KL(\rho|\alpha)+\KL(\rho|\beta)
=
\int(\sqrt{\rho_\alpha}-\sqrt{\rho_\beta})^2,
\]
when $\alpha,\beta$ have densities with respect to a common reference measure. Thus the KL marginal relaxation contains the squared Hellinger distance as its local mass-variation limit.
\end{prop}
\begin{proof}
For the upper bound, restrict to diagonal plans $\pi=(\Id,\Id)_\sharp\rho$, whose transport cost is zero and whose two marginals are both $\rho$. Conversely, after division by $\tau$, any transport cost is nonnegative and can be discarded in the lower bound, so only the best common marginal survives at the scale $\tau$. For KL, minimizing pointwise over the density $r$ of $\rho$ gives
\(r\log(r/a)-r+a+r\log(r/b)-r+b.\)
The optimality condition is $\log(r/a)+\log(r/b)=0$, hence $r=\sqrt{ab}$, and the minimum is $a+b-2\sqrt{ab}=(\sqrt a-\sqrt b)^2$.
\end{proof}
\begin{prop}[Dual of unbalanced optimal transport]
One has equality between~\eqref{eq-unbalanced-primal} and
\(\UW_c(\al,\be)=\usup{f \oplus g \leq c} - \Dd_{\psi_1}^*(-f|\al) - \Dd_{\psi_2}^*(-g|\be).\)
\end{prop}
\begin{proof}
One gets a dual by using the formula~\eqref{eq-legendre} for the dual of a divergence
\begin{align*}
\uinf{\pi \in \Mm_+(\Xx \times \Yy)} \usup{(f,g) \in \Cc(\Xx) \times \Cc(\Yy)}
\dotp{c}{\pi} + \dotp{-f}{\pi_1} + \dotp{-g}{\pi_2} - \Dd_{\psi_1}^*(-f|\al) - \Dd_{\psi_2}^*(-g|\be).
\end{align*}
Exchanging the min and the max gives
\[
\usup{(f,g)}
-\Dd_{\psi_1}^*(-f|\al)-\Dd_{\psi_2}^*(-g|\be)
+\uinf{\pi\geq0}\dotp{c-(f\oplus g)}{\pi}
=
\usup{f\oplus g\leq c}
-\Dd_{\psi_1}^*(-f|\al)-\Dd_{\psi_2}^*(-g|\be).
\]
\end{proof}
\paragraph{Reverse formulation.}
\begin{align*}
& \int_{\Xx \times \Yy} c(x,y) \d\pi(x,y)
+ \Dd_{\psi_1}(\pi_1|\al) + \Dd_{\psi_2}(\pi_2|\be) \\
=&
\int_{\Xx \times \Yy} \pa{ c(x,y) +
\psi_1\pa{ \frac{\d\pi_1}{\d\al}(x) } \frac{\d\al}{\d\pi_1}(x) +
\psi_2\pa{ \frac{\d\pi_2}{\d\be}(y) } \frac{\d\be}{\d\pi_2}(y)
} \d\pi(x,y).
\end{align*}
\eql{\label{eq-unbalanced-reverse-local-cost}
\foralls (c,r,s) \in \RR^3, \quad
L_c(r,s) \eqdef c + r\psi_1(1/r) + s\psi_2(1/s),
}
where we implicitly assume that $r\psi_1(1/r) = (\psi_1')^\infty$ when $r=0$, one has

\(\UW_c(\al,\be) = \uinf{\pi \in \Mm_+(\Xx \times \Yy)} \int L_{c(x,y)}( F(x),G(y) ) \d\pi(x,y) + \psi_1(0) \al^\bot(\Xx) + \psi_2(0) \be^\bot(\Yy)\)
\(\al = F \pi_1 + \al^\bot \qandq \be = G \pi_2 + \be^\bot.\)
Thus $F=\d\al/\d\pi_1$ and $G=\d\be/\d\pi_2$ on the absolutely continuous parts.

\paragraph{Homogeneous formulation.}
\eql{\label{eq-unbalanced-homogeneous-local-cost}
H_c(r,s) \eqdef \uinf{\th>0} \th L_c(r/\th,s/\th).
}
\eql{\label{eq-homogeneous}
\HW_c(\al,\be) = \uinf{\pi \in \Mm_+(\Xx \times \Yy)}
\int H_{c(x,y)}( F(x),G(y) ) \d\pi(x,y)
+ \psi_1(0) \al^\bot(\Xx)
+ \psi_2(0) \be^\bot(\Yy).
}
\begin{prop}[Homogenization does not change the unbalanced cost]
One has $\HW_c(\al,\be)=\UW_c(\al,\be)$.
\end{prop}
\begin{proof}
The inequality $\HW\leq\UW$ follows from $H_c\leq L_c$. For the reverse inequality, take a feasible measure $\pi$ in the homogeneous formulation. By definition of the perspective transform, for every $(x,y)$ and every $\eta>0$ there exists a scale $\theta(x,y)>0$ such that
\(H_{c(x,y)}(F(x),G(y))+\eta \geq \theta(x,y)L_{c(x,y)}\big(F(x)/\theta(x,y),G(y)/\theta(x,y)\big).\)
Replacing $\pi$ by the rescaled measure $\tilde\pi=\theta\,\pi$ and the densities by $F/\theta$ and $G/\theta$ gives an admissible competitor for the reverse formulation with cost no larger than the homogeneous cost plus $\eta\pi(\X\times\Y)$. Letting $\eta\to0$ yields $\UW\leq\HW$. The singular terms are unchanged because the same rescaling is performed before taking the Lebesgue decomposition of the marginals.
\end{proof}
\paragraph{Conic lifting.}
The last formulation is obtained by lifting the above optimization problem on the cone space $\mathfrak{C}[\Xx] \eqdef (\Xx \times \RR_+) / \sim$ associated to $\Xx$ (and similarly for $\Yy$). Here, the equivalence relation $\sim$ ``glues'' together all the points $(x,0)$, which are often called the ``apex'' of the cone.

\(\De((x,r),(y,s)) \eqdef H_{c(x,y)}(r^p,s^p)^{1/p}.\)
\begin{itemize}
\item $\Dd_\psi=\KL$, $p=2$ and $c(x,y)=-\log \cos^2(d(x,y) \wedge \frac{\pi}{2})$ in which case $\De$ is the so-called cone metric associated to $d$ and reads
\(\De((x,r),(y,s))^2 = r^2+s^2-2rs \cos( d(x,y) \wedge \frac{\pi}{2} ).\)
\item $\Dd_\psi=\KL$, $p=2$ and $c(x,y)=d(x,y)^2$, so that
\(\De((x,r),(y,s))^2 = r^2+s^2-2rs e^{-d(x,y)^2/2}.\)
\item $\Dd_\psi=\TV$, $p=1$ and $c(x,y)=d(x,y)$, so that
\(\De((x,r),(y,s)) = r+s-(r \wedge s)(2-d(x,y))_+.\)
\end{itemize}
\eq{
\CW(\al,\be) = \uinf{ \ga \in \Mm_+(\mathfrak{C}[\Xx]^2) }
\enscond{
\int \De((x,r),(y,s))^p \d\ga((x,r),(y,s))
}{
\int r^p \d \ga_1(\cdot,r) = \al, \:
\int s^p \d \ga_2(\cdot,s) = \be
}.
}
\begin{thm}[Cone formulation of unbalanced OT]\label{thm-cone-unbalanced-ot}
One has $\UW=\HW=\CW$. If $\De$ is a distance, so is $\CW^{1/p}$.
\end{thm}
\begin{proof}
The equality $\UW=\HW$ is the previous proposition. To prove $\HW=\CW$, disintegrate an admissible cone coupling $\gamma$ with respect to its spatial variables $(x,y)$ and radii $(r,s)$. The cone marginal constraints say precisely that the spatial marginals are recovered after weighting by $r^p$ and $s^p$. Because $\De((x,r),(y,s))^p=H_{c(x,y)}(r^p,s^p)$, integrating the cone cost gives the homogeneous objective. Conversely, any homogeneous competitor can be lifted to the cone by placing, over each $(x,y)$, radii whose $p$-th powers are the two density factors appearing in $H_c$.

If $\De$ is a distance on the cone, then $\CW^{1/p}$ is the usual $p$-Wasserstein distance between the lifted measures under the linear cone-marginal constraints. Symmetry and the triangle inequality follow from the corresponding Wasserstein properties and the gluing lemma on the cone. If the distance is zero, an optimal cone coupling is concentrated on the diagonal of the cone, so the weighted projections agree and therefore $\alpha=\beta$.
\end{proof}
\paragraph{Entropic regularization.}
\eq{
\uinf{\pi \in \Mm_+(\Xx \times \Yy)} \int_{\Xx \times \Yy} c(x,y) \d\pi(x,y)
+ \Dd_{\psi_1}(\pi_1|\al) + \Dd_{\psi_2}(\pi_2|\be) + \epsilon \Dd_\phi(\pi|\al \otimes \be)
}
\(\usup{f,g} - \Dd_{\psi_1}^*(-f|\al) - \Dd_{\psi_2}^*(-g|\be) - \epsilon \Dd_\phi^*\pa{ \frac{f\oplus g-c}{\epsilon} }.\)
When using $\Dd_\phi=\KL$, the primal-dual relation reads $\pi=e^{\frac{f\oplus g - c}{\epsilon}}$ and when also considering $\Dd_{\psi_1}=\Dd_{\psi_2} = \tau \KL$, the dual reads

\eq{
\usup{f,g}
- \tau \pa{ \dotp{e^{\frac{-f}{\tau}}}{\al} + \dotp{e^{\frac{-g}{\tau}}}{\be} - 2 }
- \epsilon \iint_{\Xx \times \Yy} e^{\frac{f\oplus g - c}{\epsilon}} \d \al \otimes \d \be + \epsilon.
}
\begin{align*}
f &= -\frac{\tau\epsilon}{\tau+\epsilon} \log\int_\Yy \exp\pa{ \frac{g(y) - c(\cdot,y)}{\epsilon} } \d\be(y), \\
g &= -\frac{\tau\epsilon}{\tau+\epsilon} \log\int_\Xx \exp\pa{ \frac{f(x) - c(x,\cdot)}{\epsilon} } \d\al(x). \\
\end{align*}
\subsection{OT Barycenters}
\label{sec-barycenters}
\paragraph{Fr\'echet means.}
For discrete input histograms $\{\b_s\}_{s=1}^S$, with $\b_s \in \simplex_{n_s}$, and weights $\la \in \simplex_S$, a Wasserstein barycenter can be computed by minimizing

\eql{\label{eq-wass-discr}
\umin{\a \in \simplex_n} \sum_{s=1}^S \la_s \MKD_{\C_s}(\a,\b_s),
}
where the cost matrices $\C_s \in \RR^{n \times n_s}$ are prescribed.

Given a set of input measures $(\be_s)_s$ defined on some space $\X$, the barycenter problem becomes

\eql{\label{eq-barycenter-generic}
\umin{\al \in \Mm_+^1(\X)} \sum_{s=1}^S \la_s \MK_{\c}(\al,\be_s).
}
\begin{rem}[Two measures give a Wasserstein geodesic]
\end{rem}
\begin{example}[Dirac inputs]
Problem~\eqref{eq-barycenter-generic} generalizes the computation of barycenters of points $(x_s)_{s=1}^S \in \X^S$ to arbitrary measures. Indeed, if $\be_s=\de_{x_s}$ is a single Dirac mass, then a solution to~\eqref{eq-barycenter-generic} is $\de_{x^\star}$ where $x^\star$ is a Fr\'echet mean of the points $(x_s)_s$.

\end{example}
\begin{rem}[Mean of a quadratic barycenter]
For $c(x,y)=\norm{x-y}^2$, the mean of the barycenter $\al^\star$ is necessarily the barycenter of the means,

\(\int_\Xx x \d\al^\star(x) = \sum_s \la_s \int_\Xx x \d\be_s(x).\)
\end{rem}
\begin{proposition}[Convexity of the OT cost]
$(\al,\be) \mapsto \MK_{\c}(\al,\be)$ is convex.
\end{proposition}
\begin{proof}
Let $(\al_0,\be_0)$ and $(\al_1,\be_1)$ be two pairs of probability measures and let $t \in [0,1]$. For $\epsilon>0$, choose couplings $\pi_i \in \Couplings(\al_i,\be_i)$ such that
\(\int c \d\pi_i \leq \MK_\c(\al_i,\be_i)+\epsilon \qforq i=0,1.\)
Then $\pi_t \eqdef (1-t)\pi_0+t\pi_1$ is a coupling between $(1-t)\al_0+t\al_1$ and $(1-t)\be_0+t\be_1$. Hence
\(\MK_\c((1-t)\al_0+t\al_1,(1-t)\be_0+t\be_1) \leq (1-t)\MK_\c(\al_0,\be_0)+t\MK_\c(\al_1,\be_1)+\epsilon.\)
Letting $\epsilon \rightarrow 0$ gives the claim.
\end{proof}
\paragraph{1-D Case.}
\begin{prop}[Quantile barycenters on the line]
For $\X=\RR$ and $c(x,y)=|x-y|^2$, the quantile function of a Wasserstein barycenter is the weighted average of the input quantile functions:
\(\cumul{\al^\star}^{-1}(r) = \sum_{s=1}^S \la_s \cumul{\be_s}^{-1}(r) \qforq r \in [0,1].\)
\end{prop}
\begin{proof}
The one-dimensional formula~\eqref{eq-wass-cumul} gives \(\sum_s \la_s \Wass_2^2(\al,\be_s) = \int_0^1 \sum_s \la_s \abs{\cumul{\al}^{-1}(r)-\cumul{\be_s}^{-1}(r)}^2 \d r.\)
The minimization decouples pointwise in $r$. For each fixed $r$, the minimizer of
$z \mapsto \sum_s \la_s |z-\cumul{\be_s}^{-1}(r)|^2$
is the weighted average $\sum_s \la_s \cumul{\be_s}^{-1}(r)$. This function is nondecreasing because it is a positive weighted sum of nondecreasing quantile functions, hence it is itself a valid quantile function.
\end{proof}
\paragraph{Gaussian Case.}
\begin{rem}[Gaussian barycenters]
\(\cov \mapsto \sum_s \la_s \Bb(\cov,\cov_s)^2,\)
\(\cov = \sum_s \la_s \pa{\cov^{1/2}\cov_s\cov^{1/2}}^{1/2}.\)
\end{rem}
\paragraph{Sinkhorn for barycenters.}
\eql{\label{eq-entropic-bary}
\umin{\a \in \simplex_n} \sum_{s=1}^S \la_s \MKD_{\C_s}^\epsilon(\a,\b_s)
}
\eql{\label{eq-bary-entropy-couplings}
\umin{ (\P_s)_s } \enscond{ \sum_{s} \la_s \KLD( \P_s|\K_s ) }{
\foralls s, \transp{\P_s}\ones_n = \b_s, \:
\P_1\ones_{n_1} = \ldots = \P_S\ones_{n_S}
}
}
where we denoted $\K_s \eqdef e^{-\C_s/\epsilon}$. Here, the barycenter $\a$ is implicitly encoded in the row marginals of all the couplings $\P_s \in \RR^{n \times n_s}$ as $\a = \P_1\ones_{n_1} = \ldots = \P_S\ones_{n_S}$.

\eql{\label{eq-bary-opt}
\P_s=\diag(\uD_s)\K_s\diag(\vD_s),
}
\begin{align}\label{eq-sinkhorn-bary}
\foralls s \in \range{1,S}, \quad \itt{\vD}_s &\eqdef \frac{\b_s}{\transp{\K_s} \it{\uD}_s}, \\
\foralls s \in \range{1,S}, \quad \itt{\uD}_s &\eqdef \frac{\itt{\a}}{\K_s \itt{\vD}_s}, \label{eq-sinkhorn-bary-2}\\
\qwhereq
\itt{\a} &\eqdef \prod_s ( \K_s \itt{\vD}_s )^{\la_s}. \label{eq-sinkhorn-bary-3}
\end{align}
\begin{prop}[Dual of entropic barycenters]
The optimal $(\uD_s,\vD_s)$ appearing in~\eqref{eq-bary-opt} can be written as $(\uD_s,\vD_s) = ( e^{\fD_s/\epsilon},e^{\gD_s/\epsilon} )$ where $( \fD_s,\gD_s )_s$ are the solutions of the following program (whose value matches the one of \eqref{eq-entropic-bary})
\eql{\label{eq-dual-bary-entropy}
\umax{ ( \fD_s,\gD_s )_s } \enscond{
\sum_s \la_s \pa{
\dotp{\gD_s}{\b_s} - \epsilon \dotp{\K_s e^{\gD_s/\epsilon}}{e^{\fD_s/\epsilon}}
}
}{
\sum_s \la_s \fD_s = 0
}.
}
\end{prop}
\begin{proof}
Introducing Lagrange multipliers in~\eqref{eq-bary-entropy-couplings} leads to
\begin{align*}
\umin{ (\P_s)_s, \a }
\umax{ ( \fD_s,\gD_s )_s }
\sum_{s} \la_s \Big(
\epsilon \KLD( \P_s|\K_s )
+
\dotp{\a - \P_s\ones_{n_s}}{ \fD_s } \\
\qquad\qquad	+
\dotp{\b_s - \transp{\P_s}\ones_n}{ \gD_s }
\Big).
\end{align*}
Strong duality holds, so one can exchange the minimum and the maximum and obtain
\begin{align*}
&\umax{ ( \fD_s,\gD_s )_s }
\sum_s \la_s \pa{
\dotp{\gD_s}{\b_s}
+
\umin{\P_s} \epsilon \KLD( \P_s|\K_s ) - \dotp{\P_s}{\fD_s \oplus \gD_s }
} \\
&\qquad\qquad	+ \umin{\a}
\dotp{ \sum_{s} \la_s \fD_s }{\a}.
\end{align*}
The explicit minimization on $\a$ gives the constraint $\sum_{s} \la_s \fD_s=0$ together with
\begin{align*}
\umax{ ( \fD_s,\gD_s )_s }
\sum_s \la_s
\dotp{\gD_s}{\b_s}
-
\epsilon \KLD^*\pa{ \frac{\fD_s \oplus \gD_s}{\epsilon}|\K_s }
\end{align*}
where $\KLD^*(\cdot|\K_s)$ is the Legendre transform~\eqref{eq-legendre} of the function $\KLD(\cdot|\K_s)$.

This Legendre transform reads
\eql{\label{eq-legendre-kl}
\KLD^*(\VectMode{U}|\K) = \sum_{i,j} \K_{i,j} (e^{\VectMode{U}_{i,j}}-1),
}
which shows the desired formula.

To show~\eqref{eq-legendre-kl}, since this function is separable, one needs to compute
\(\foralls (u,k) \in \RR \times \RR_+^*, \quad \KLD^*(u|k) \eqdef \umax{r \geq 0} ur - \pa{ r \log(r/k) - r + k }\)
whose optimality condition reads $u=\log(r/k)$, i.e. $r = k e^{u}$, hence the result. The case $k=0$ follows by lower semicontinuity and gives zero.
\end{proof}
\subsection{Multimarginal OT}
\label{subsec:multi-marginal}
\paragraph{Definition, example and properties.}
The multi-marginal formulation replaces a coupling between two measures by a joint distribution with several prescribed marginals. Given measures $(\al_s)_{s=1}^S$ on spaces $(\X_s)_{s=1}^S$ and a cost $c: \X_1 \times \cdots \times \X_S \to \RR$, the problem reads

\(\umin{\pi \in \Couplings(\al_1,\ldots,\al_S)} \int_{\X_1 \times \cdots \times \X_S} c(x_1,\ldots,x_S)\d\pi(x_1,\ldots,x_S),\)
where $\Couplings(\al_1,\ldots,\al_S)$ is the set of probability measures whose $s$-th marginal is $\al_s$. This is still a linear program in the discrete setting, but its ambient dimension grows exponentially with the number of marginals.

\paragraph{Multi-marginal formulation of barycenters.}
\(c(x_1,\ldots,x_S) = \umin{x \in \RR^d} \sum_{s=1}^S \la_s \norm{x-x_s}^2.\)
\begin{prop}[Multi-marginal formula for quadratic barycenters]\label{prop-multimarginal-barycenter}
Let $\be_1,\ldots,\be_S\in\Pp_2(\RR^d)$ and $\la\in\simplex_S$. Define
\(B(x_1,\ldots,x_S)=\sum_{s=1}^S\la_s x_s, \qquad c_{\mathrm{bar}}(x_1,\ldots,x_S)=\min_x\sum_s\la_s\norm{x-x_s}^2.\)
If $\pi^\star$ solves the multi-marginal OT problem with marginals $(\be_s)_s$ and cost $c_{\mathrm{bar}}$, then $\al^\star=B_\sharp\pi^\star$ is a Wasserstein barycenter. Conversely, every barycenter is obtained this way from an optimal multi-marginal plan.
\end{prop}
\begin{proof}
For any candidate barycenter $\al$ and couplings $\pi_s\in\Couplings(\al,\be_s)$, glue the couplings along their common $\al$ marginal to obtain a joint law of $(X,Y_1,\ldots,Y_S)$. Conditioning on $(Y_s)_s$ and minimizing over $X$ gives
\(\sum_s\la_s\EE\norm{X-Y_s}^2 \geq \EE\min_x\sum_s\la_s\norm{x-Y_s}^2 = \EE c_{\mathrm{bar}}(Y_1,\ldots,Y_S).\)
Taking the infimum over the couplings gives that the barycenter value is at least the multi-marginal value. Conversely, from an optimal multi-marginal plan $\pi^\star$, set $X=B(Y_1,\ldots,Y_S)$. The couplings between $X$ and each $Y_s$ are feasible for the barycenter problem and attain exactly the multi-marginal cost, proving equality and the formula.
\end{proof}
\begin{cor}[Gaussian and discrete barycenters]\label{cor-gaussian-discrete-barycenters}
Quadratic Wasserstein barycenters of Gaussian measures are Gaussian. If the input measures are discrete, then there exists a barycenter supported on the set of weighted averages $\sum_s\la_s x_{s,i_s}$ of one support point from each input; in particular, if the $s$-th input has $n_s$ atoms, a barycenter exists with at most $\prod_s n_s$ atoms.
\end{cor}
\begin{proof}
Let the input Gaussians have means $\mean_s$ and covariances $\cov_s$. For any candidate barycenter $\al$ with mean $\mean$ and covariance $\cov$, Gelbrich's inequality gives
\(\Wass_2^2(\al,\be_s) \geq \norm{\mean-\mean_s}^2+\Bb(\cov,\cov_s)^2,\)
with equality when $\al$ is Gaussian. Therefore the barycenter objective is bounded below by a function depending only on $(\mean,\cov)$, and this lower bound is attained by the Gaussian measure with the minimizing mean and covariance. Hence at least one barycenter is Gaussian, and uniqueness in the usual nondegenerate setting gives the Gaussian barycenter mentioned above. For discrete inputs, any multi-marginal optimizer is supported on the finite product of the input supports, and $B$ maps this product to at most $\prod_s n_s$ points.
\end{proof}
\paragraph{Entropic regularization of multi-marginal.}
\(\inf_{\pi\in\Couplings(\al_1,\ldots,\al_S)} \int c\d\pi+\epsilon\KL(\pi|\al_1\otimes\cdots\otimes\al_S).\)
\[
\d\pi^\star(x_1,\ldots,x_S)
=
\exp\!\left(\frac{\sum_s f_s(x_s)-c(x_1,\ldots,x_S)}{\epsilon}\right)
\prod_s\d\al_s(x_s),
\]
\subsection{Sliced Wasserstein}
\label{sec-sliced-wasserstein}
For measures on $\RR^d$, sliced Wasserstein distances replace one high-dimensional transport problem by many one-dimensional ones. For $\theta \in \Sphere^{d-1}$, let $P_\theta(x)=\dotp{x}{\theta}$ be the projection on direction $\theta$. The sliced distance is

\eql{\label{eq-sliced-wasserstein}
\SW_p^p(\al,\be)
\eqdef
\int_{\Sphere^{d-1}}
\Wass_p^p\pa{
(P_\theta)_\sharp \al,
(P_\theta)_\sharp \be
}\d\theta.
}
\begin{prop}[Metric properties of sliced Wasserstein]\label{prop-sliced-wasserstein-metric}
For $p\geq1$, $\SW_p$ is a distance on $\Pp_p(\RR^d)$. Moreover, $\SW_p$ metrizes weak convergence together with convergence of the $p$-th moment. Finally, \(\SW_p(\alpha,\beta)\leq \Wass_p(\alpha,\beta),\)
and, for $p=2$ with the uniform probability measure on the sphere, \(\SW_2^2(\alpha,\beta)\leq \frac1d \Wass_2^2(\alpha,\beta).\)
\end{prop}
\begin{proof}
Non-negativity and symmetry follow from the one-dimensional Wasserstein distance. For the triangle inequality, apply the triangle inequality of $\Wass_p$ for each direction $\theta$ and then Minkowski's inequality in $L^p(\Sphere^{d-1})$.

If $\SW_p(\alpha,\beta)=0$, then $(P_\theta)_\sharp\alpha=(P_\theta)_\sharp\beta$ for almost every direction, hence for all directions by continuity of characteristic functions. The Cramer--Wold theorem then implies $\alpha=\beta$. This proves that $\SW_p$ is a distance.

The bound $\SW_p\leq\Wass_p$ follows because $P_\theta$ is $1$-Lipschitz, so
\(\Wass_p((P_\theta)_\sharp\alpha,(P_\theta)_\sharp\beta) \leq \Wass_p(\alpha,\beta)\)
for every $\theta$. For $p=2$, using any coupling $\pi$ between $\alpha$ and $\beta$, \(\int_{\Sphere^{d-1}}\int |\dotp{x-y}{\theta}|^2\d\pi(x,y)\d\theta = \frac1d\int\norm{x-y}^2\d\pi(x,y).\)
Optimizing over $\pi$ gives the sharper displayed inequality.

The weak-convergence statement follows from the same Cramer--Wold mechanism plus the moment condition: convergence in $\SW_p$ gives convergence of almost all one-dimensional projections and tightness of the $p$-moments; conversely, weak convergence with $p$-moment convergence implies convergence of projected $\Wass_p$ distances and dominated convergence on the sphere. Detailed proofs and statistical refinements can be found in.
\end{proof}
\begin{rem}[Hilbert embedding for $\SW_2$]
In one dimension, $\Wass_2$ is the $L^2(0,1)$ distance between quantile functions. Hence

\(\SW_2^2(\alpha,\beta) = \int_{\Sphere^{d-1}} \int_0^1 \abs{F^{-1}_{\theta,\alpha}(u)-F^{-1}_{\theta,\beta}(u)}^2 \d u\,\d\theta,\)
where $F^{-1}_{\theta,\alpha}$ is the quantile of $(P_\theta)_\sharp\alpha$. Thus $\SW_2$ is a Hilbertian distance after embedding each measure into its field of projected quantiles.

\end{rem}
\begin{defn}[Max-, min- and subspace-sliced variants]\label{def-sliced-variants}
The max-sliced distance replaces the average over directions by the supremum
\(\MaxSW_p(\alpha,\beta)= \sup_{\theta\in\Sphere^{d-1}} \Wass_p((P_\theta)_\sharp\alpha,(P_\theta)_\sharp\beta).\)
Min-sliced transport-plan variants such as min-SWGG instead choose a direction whose lifted 1-D matching has minimal full-dimensional cost; they provide a cheap transport plan and an upper bound on $\Wass_2^2$ rather than a lower-dimensional distance. Subspace-sliced variants replace lines by $k$-dimensional orthogonal projectors $U^\top x$, with $U^\top U=I_k$. Max-sliced variants are useful in generative modeling.
\end{defn}
\begin{prop}[Basic bounds for sliced variants]\label{prop-sliced-variant-bounds}
With normalized surface measure, \(\SW_p(\alpha,\beta) \leq \MaxSW_p(\alpha,\beta) \leq \Wass_p(\alpha,\beta).\)
For min-SWGG-type lifted plans, the induced cost is an upper bound on $\Wass_2^2(\alpha,\beta)$. For $k$-dimensional subspace projections, the corresponding projected Wasserstein distance is bounded above by $\Wass_p$ for every fixed subspace.
\end{prop}
\begin{proof}
The first inequality follows because an average is bounded by a supremum. The second and the subspace bound follow because orthogonal projections are $1$-Lipschitz, so pushing any admissible coupling forward gives a feasible coupling for the projected problem with no larger cost. A lifted one-dimensional matching is an admissible coupling between the original measures; evaluating the full quadratic cost on any admissible coupling can only be larger than the optimal $\Wass_2^2$ value.
\end{proof}
\begin{defn}[Subspace robust Wasserstein]\label{def-subspace-robust-wasserstein}
For a coupling $\pi$ between measures on $\RR^d$, define its displacement covariance
\(M_\pi=\int (x-y)(x-y)^\top\d\pi(x,y).\)
The subspace robust Wasserstein cost of rank $k$ is \(\SRW_{2,k}^2(\alpha,\beta) \eqdef \inf_{\pi\in\Couplings(\alpha,\beta)} \sum_{\ell=1}^k \lambda_\ell(M_\pi),\)
where $\lambda_1\geq\cdots\geq\lambda_d$ are eigenvalues. Taking the trace, i.e. summing all eigenvalues, recovers $\Wass_2^2$.
\end{defn}
\begin{prop}[Metric behavior of robust subspace costs]\label{prop-subspace-robust-metric}
The relaxed covariance-gauge formulation above is a convex minimization over couplings; replacing the trace gauge by the spectral partial-sum gauge emphasizes the largest displacement directions and satisfies $\SRW_{2,k}\leq\Wass_2$ when the eigenvalue sum is normalized consistently.
\end{prop}
\begin{proof}
Metricity of the projected max formulation follows from applying the Wasserstein metric after each orthogonal projection and taking a supremum over projections that separates points. The bound by $\Wass_2$ again follows from $1$-Lipschitzness of projections. For the relaxed formulation, $M_\pi$ depends linearly on $\pi$, and the sum of the largest $k$ eigenvalues is a convex spectral gauge; minimizing a convex function over the convex coupling set gives a convex problem. Since the trace is the sum of all eigenvalues and dominates the normalized partial sum, the value is bounded by the usual quadratic OT cost. This is the finite-dimensional robust-subspace viewpoint; related spectral Wasserstein gauges give analogous constructions.
\end{proof}
\subsection{Linear Optimal Transport}
\label{sec-linear-ot}
Fix a reference absolutely continuous measure $\rho$ and let $T_\alpha$ be the Brenier map pushing $\rho$ to $\alpha$. The linear OT embedding is

\begin{equation}\label{eq-lot-embedding}
\alpha \mapsto T_\alpha-\Id\in L^2(\rho;\RR^d),
\qquad
\LOT_\rho(\alpha,\beta)=\norm{T_\alpha-T_\beta}_{L^2(\rho)}.
\end{equation}
\begin{prop}[Local stability of linear OT]\label{prop-linear-ot-stability}
Assume that the measures are supported on a fixed convex compact set, with densities bounded above and below, and that the Brenier maps from $\rho$ are regular. Then locally around $\rho$,
\(\Wass_2(\alpha,\beta)\leq \LOT_\rho(\alpha,\beta) \quad\text{and}\quad \LOT_\rho(\alpha,\beta)\leq C\Wass_2(\alpha,\beta)^\eta\)
for constants $C>0$ and $\eta\in(0,1]$ depending on regularity.
\end{prop}
\begin{proof}
The first inequality is immediate: $(T_\alpha,T_\beta)_\sharp\rho$ is a feasible coupling between $\alpha$ and $\beta$. The reverse local estimate is a standard stability statement for the Monge--Amp\`ere equation under the stated regularity assumptions: changes in the target measure control changes in the Brenier potential in H\"older norms, hence control $T_\alpha-T_\beta$ in $L^2(\rho)$. In simple one-dimensional settings, quantile functions make this exact with $\eta=1$.
\end{proof}
\subsection{Weak Optimal Transport}
Given $\pi\in\Couplings(\alpha,\beta)$, disintegrate it as $\pi(\d x,\d y)=\pi_x(\d y)\alpha(\d x)$. A weak OT problem has the form

\(\inf_{\pi\in\Couplings(\alpha,\beta)} \int C(x,\pi_x)\d\alpha(x),\)
where $C(x,\cdot)$ is convex on probability measures. The classical Kantorovich problem is recovered when $C(x,\nu)=\int c(x,y)\d\nu(y)$.

\subsection{Gromov--Wasserstein}
\label{sec-gromov-wasserstein}
\begin{prop}[Euclidean GW is controlled by Wasserstein]\label{prop-gw-controlled-by-wasserstein}
Let $\alpha,\beta$ be probability measures on $\RR^d$, equipped with the Euclidean distance, and take $\De(u,v)=|u-v|$ in~\eqref{eq-gw-generic}. Then
\(\GW((\RR^d,\norm{\cdot},\alpha),(\RR^d,\norm{\cdot},\beta)) \leq 2\Wass_p(\alpha,\beta).\)
\end{prop}
\begin{proof}
Let $\pi$ be any coupling between $\alpha$ and $\beta$. For two independent pairs $(X,Y),(X',Y')\sim\pi$, the reverse triangle inequality gives \(\big|\norm{X-X'}-\norm{Y-Y'}\big| \leq \norm{X-Y}+\norm{X'-Y'}.\)
Taking the $L^p$ norm and using Minkowski gives a bound by $2(\int\norm{x-y}^p\d\pi)^{1/p}$. Optimizing over $\pi$ proves the claim.
\end{proof}
Optimal transport needs a ground cost $\C$ to compare histograms $(\a,\b)$, and thus cannot be used directly if the histograms are not defined on the same underlying space, or if one cannot pre-register these spaces to define a ground cost.

To address this issue, one can instead use a weaker requirement: two matrices $\distD \in \RR^{n \times n}$ and $\distD' \in \RR^{m \times m}$ are available and represent relationships between the points on which the histograms are defined. A typical scenario is when these matrices are powers of distance matrices.

\eql{\label{eq-gw-def}
\GWD( (\a,\distD), (\b,\distD') )^p \eqdef
\umin{ \P \in \CouplingsD(\a,\b) }
\Ee_{\distD,\distD'}(\P) \eqdef
\sum_{i,j,i',j'} \De(\distD_{i,i'},\distD'_{j,j'})^p \P_{i,j}\P_{i',j'},
}
where $p \geq 1$ and $\De$ is a distance on $\RR$, typically $\De(u,v)=|u-v|$.

\paragraph{General setting.}
The general setting corresponds to computing couplings between metric measure spaces $\XX=(\X,\dist_\X,\al)$ and $\YY=(\Y,\dist_\Y,\be)$ where $(\dist_\X,\dist_\Y)$ are distances and $(\al,\be)$ are measures on their respective spaces.

\begin{align}
\label{eq-gw-generic}
\GW( \XX,\YY )^p \eqdef
\umin{ \pi \in \Couplings(\al,\be) }
\int_{\X^2 \times \Y^2}
\De(\dist_\X(x,x'),\dist_\Y(y,y'))^p
\d\pi(x,y)\d\pi(x',y').
\end{align}
\begin{thm}[Gromov--Wasserstein metric modulo isometries]
$\GW$ defines a distance between metric measure spaces up to isometries.
\end{thm}
\begin{proof}
If $\GW( \XX,\YY )=0$ and $\pi$ is an optimal plan, then $\dist_\X(x,x')=\dist_\Y(y,y')$ holds $\pi \otimes \pi$-almost everywhere. By continuity, this equality holds on $\supp(\pi)^2$.

\(\dist_\pi((x,y),(x',y')) \eqdef \frac{1}{2} \dist_\X(x,x') + \frac{1}{2} \dist_\Y(y,y').\)

One has, for $((x,y),(x',y')) \in \supp(\pi)^2$, $\dist_\X(x,x')=\dist_\Y(y,y')$ and thus
\(\dist_\X(\psi(x,y),\psi(x',y')) = \dist_\X(x,x') = \dist_\Y(y,y') = \dist_\pi((x,y),(x',y')),\)
thus $\psi$ is an isometry, hence injective. Indeed, since $\psi_\sharp \pi = \al$, there exists a sequence $(x_k,y_k) \in \supp(\pi)$ such that $x_k \rightarrow x$.

Since $\dist_\X(x_k,x_{k'})=\dist_\Y(y_k,y_{k'})$ the sequence $(y_k)_k$ is Cauchy, and since $\supp(\pi)$ is closed, it is complete and thus $(x_k,y_k)$ is converging to some $(x,y) \in \supp(\pi)$. The same argument for the second projection shows that this support space is also isometric to $\YY$.

Given an optimal $\pi$ and $\xi$ between three spaces $(\Xx,\Yy,\Zz)$ with measures $(\al,\be,\ga)$, we use the gluing Lemma~\ref{lem-gluing-general} to define $\si$ with marginals $(P_{\Xx,\Yy})_\sharp \si = \pi$ and $(P_{\Yy,\Zz})_\sharp \si = \xi$.
We denote $\rho \eqdef (P_{\Xx,\Zz})_\sharp \si$ the composition of the two couplings, so that
$$\begin{aligned}
\GW( \XX,\ZZ ) & \leq \Big( \int_{\Xx \times \Zz} \De(\dist_\Xx(x,x'),\dist_\Zz(z,z'))^p \d\rho(x,z)\d\rho(x',z') \Big)^{1/p}
\\
&\leq \Big( \int_{\Xx \times \Yy \times \Zz} ( \De(\dist_\Xx(x,x'),\dist_\Yy(y,y')) + \De(\dist_\Yy(y,y'),\dist_\Zz(z,z')) )^p \d\si(x,y,z) \d\si(x',y',z')\Big)^{1/p} \\
&\leq \Big( \int_{\Xx \times \Yy \times \Zz} \De(\dist_\Xx(x,x'),\dist_\Yy(y,y'))^p \d\si(x,y,z) \d\si(x',y',z')\Big)^{1/p} \\
&\qquad +
\Big( \int_{\Xx \times \Yy \times \Zz} \De(\dist_\Yy(y,y'),\dist_\Zz(z,z'))^p \d\si(x,y,z) \d\si(x',y',z')\Big)^{1/p}
\\
&\leq \GW( \XX,\YY ) + \GW( \YY,\ZZ ),
\end{aligned}$$
where the first inequality follows from the feasibility of $\rho$, the second is the triangular inequality of $\De$,
and the third is the Minkowski inequality.
\end{proof}
\begin{rem}[Gromov--Wasserstein geodesics]
\(\dist_t((x_0,x_1), (x_0',x_1')) \eqdef (1-t)\dist_{\X_0}(x_0,x_0') + t\dist_{\X_1}(x_1,x_1').\)
\end{rem}
\paragraph{Entropic regularization and iterative solver.}
\eql{\label{eq-gw-entropy}
\umin{ \P \in \CouplingsD(\a,\b) }
\Ee_{\distD,\distD'}(\P) - \epsilon \HD(\P).
}
\eql{\label{eq-gw-sinkh}
\itt{\P} \eqdef \umin{ \P \in \CouplingsD(\a,\b) } \dotp{\P}{\it{\C}} - \epsilon\H(\P)
\qwhereq
}
\(\it{\C} \eqdef \distD^{\odot 2}\a\,\ones_m^\top + \ones_n(\distD'^{\odot 2}\b)^\top - 2\distD\,\it{\P}\,\transp{\distD'},\)
\subsection{Metric learning and inverse OT}
\paragraph{Metric learning and derivatives of OT.}
\begin{prop}[Derivative with respect to the cost]
In the discrete setting, assume that the optimal coupling for $\MKD_\C(\a,\b)$ is unique and denote it by $\P^\star(\C)$. Then $\C \mapsto \MKD_\C(\a,\b)$ is differentiable at $\C$ and
\(\nabla_\C \MKD_\C(\a,\b)=\P^\star(\C).\)
\end{prop}
\begin{proof}
The value is the minimum of affine functions of $\C$, \(\MKD_\C(\a,\b)=\min_{\P\in\CouplingsD(\a,\b)}\dotp{\C}{\P}.\)
The envelope theorem, or equivalently Danskin's theorem, states that the subdifferential with respect to $\C$ is the convex hull of the optimal couplings. If the optimizer is unique, this subdifferential is the singleton $\{\P^\star(\C)\}$, hence the value is differentiable with the displayed gradient.
\end{proof}
Thus, if the cost is parameterized as $\C_\theta$, gradients of losses involving OT values are obtained by backpropagating through the inner product $\dotp{\P^\star(\C_\theta)}{\partial_\theta \C_\theta}$. The difficulty is not differentiating a solved OT problem, but learning a cost for which the resulting matching has the desired semantic behavior; this is a bilevel and usually non-convex optimization problem.

\paragraph{Inverse Optimal Transport.}
\(\inf_{\pi\in\Couplings(\alpha,\beta)}\int c(x,y)\d\pi(x,y).\)
A useful statistical methodology is to measure the suboptimality of the observed plan through a Fenchel--Young loss. Write the score as $s=-c$ and define the convex regularized prediction value

\(G_\epsilon(s)= \sup_{\pi\in\Couplings(\alpha,\beta)} \int s\d\pi-\epsilon\KL(\pi|\alpha\otimes\beta).\)
\(\mathcal L_\epsilon(c;\widehat\pi) = G_\epsilon(-c)+G_\epsilon^*(\widehat\pi)+\int c\d\widehat\pi\)
is nonnegative by Fenchel's inequality and vanishes exactly when $\widehat\pi\in\partial G_\epsilon(-c)$, i.e. when $\widehat\pi$ satisfies the regularized optimality conditions for $c$. Entropic regularization is important here because it makes the forward map smoother and provides gradients with respect to $c$, at the price of a bias that must be controlled statistically.

\(C_\theta=\sum_{r=1}^R \theta_r C^{(r)}, \qquad \theta\in\Theta,\)
where $\Theta$ is convex and the matrices $C^{(r)}$ encode features, graph distances or a Mahalanobis parameterization.

\begin{prop}[Convex dual-gap formulation of inverse OT]\label{prop-inverse-ot-convex}
Let $\widehat P\in\CouplingsD(a,b)$ be an observed coupling and let $C_\theta$ depend affinely on $\theta\in\Theta$, where $\Theta$ is convex. The condition that $\widehat P$ is optimal for the cost $C_\theta$ is equivalent to the existence of dual potentials $(f,g)$ such that
\(f_i+g_j\leq (C_\theta)_{i,j} \qandq \sum_{i,j}\widehat P_{i,j}\big((C_\theta)_{i,j}-f_i-g_j\big)=0.\)
\eql{\label{eq-inverse-ot-convex}
\umin{\theta\in\Theta,f,g}
R(\theta)+\lambda
\sum_{i,j}\widehat P_{i,j}\big((C_\theta)_{i,j}-f_i-g_j\big)
\quad\text{subject to}\quad
f_i+g_j\leq (C_\theta)_{i,j}\quad\forall i,j.
}
\end{prop}
\begin{proof}
For a fixed cost $C_\theta$, Kantorovich duality gives \(\min_{P\in\CouplingsD(a,b)}\dotp{C_\theta}{P} = \max_{f_i+g_j\leq (C_\theta)_{i,j}} \dotp{f}{a}+\dotp{g}{b}.\)
Since $\widehat P$ has marginals $(a,b)$, every dual feasible pair satisfies \(\dotp{C_\theta}{\widehat P}-\dotp{f}{a}-\dotp{g}{b} = \sum_{i,j}\widehat P_{i,j}\big((C_\theta)_{i,j}-f_i-g_j\big) \geq0.\)
It vanishes if and only if $\widehat P$ reaches the dual value and is therefore optimal. If $C_\theta$ is affine and $\Theta$ and $R$ are convex, the constraints and objective in~\eqref{eq-inverse-ot-convex} are convex, proving the relaxation claim.
\end{proof}
\(\widehat P_{i,j}\approx a_i b_j \exp\pa{\frac{f_i+g_j-(C_\theta)_{i,j}}{\epsilon}},\)